\subsection{What the Democratic Dividend Leaves Out}
\label{sec:8.7.7}
Section 8.7.6 credits a state with functioning rule-of-law institutions, a free press and usable courts with a democratic dividend: the same technology, deployed with checks on it. Stated that way the dividend is a property a state either has or lacks, and I no longer think that is right. Having those institutions is necessary and it is not sufficient, and the ways it fails to be sufficient are separable. Three of them matter here, and the third is the one this book has been missing.

\runin{First: reach}
A democracy's AI can be fully accountable to its electorate and entirely unaccountable to the people it is pointed at. Every mechanism in the dividend runs through the franchise — you vote out the government that authorized the system, you read the press that covers it, you sue in courts your citizenship gives you standing in. An occupied population, a foreign population, a population on the other side of a border has access to none of that. The check reaches exactly as far as the demos and stops. A system built by a democracy and aimed outward is, from the perspective of the people it selects, indistinguishable in its accountability from one built by a state with no elections at all.

\runin{Second: findings without consequences}
The mechanisms can work as mechanisms and fail as checks. In the Israeli case the apparatus performed: journalists published the targeting systems in detail, serving personnel spoke to them, and the courts were simultaneously trying a sitting prime minister. Every institution did what the dividend says institutions do. No finding attached to a body with standing to stop what it described.

The American case shows the same decoupling in a different register. A civilian-harm investigation into a strike that killed at least a hundred and sixty-five people, most of them schoolchildren, was completed and not released; the lever available to the Senate committee that wanted it was a provision in the annual defense authorization withholding three-quarters of the Defense Secretary's travel budget until the unredacted investigation was handed over \autocite{stassis2026senate}. That is a co-equal branch of government reduced to expense-account leverage over a war-crimes investigation. The oversight mechanism existed, was used, and produced a travel-budget dispute. This is the same decoupling section 8.6.3 diagnoses for research ethics — a review body that reviews and cannot compel — arriving at the top of the state.

\runin{Third: the countermajoritarian half}
The deepest problem is that the dividend, as this book has been describing it, is entirely majoritarian. Every mechanism in it translates public will into outcome: elections aggregate preference, a free press informs the aggregation, and the implied failure mode is poor transmission — the public wanted something better and the institutions failed to deliver it.

But if the consensus is itself the problem, better transmission makes things worse. A well-instrumented democracy pointed at an unpopular minority is a democracy that finds out efficiently what the majority wants done and does it. Nothing in an aggregation mechanism protects anyone from the aggregate.

What is missing is the machinery that constrains preference instead of aggregating it: rights that hold regardless of being popular, courts that rule against a government which won the last election fairly, constitutional limits not repealable by whoever holds a majority this year, and standing for people outside the demos entirely — which is the whole function of international humanitarian law, a body of rules whose entire purpose is to bind a state toward people who will never vote in its elections. That is not an argument against democracy. It is the half of democracy this book has left out, and leaving it out is what made the dividend look like a property of having elections.

Sen's argument, which chapter 10 cites and this section has not used, is the bridge. What makes democracy valuable in his account is public reasoning and the capacity to demand accountability, and neither is the ballot. Both survive the majority being wrong; the ballot does not.

This is what finally lets the book say what makes a system antifascist rather than merely responsive. It is not that the system tracks what people want. It is that there are things it will not do to a person no matter who wants it. Chapter 3 sets out what such a floor requires, and section 3.1 what each way of building one costs. A system with no such floor is not protected by being deployed in a democracy — it is a system that will faithfully execute whatever the demos decides, which is exactly the capability an authoritarian movement that wins an election requires.
